\section{Experimental Setup}
\label{sec:experimental-setup}

\subsection{Dataset and Splits}
All reported numbers in this report use the official \dataset validation split. That split contains 10{,}895 natural-language queries over 2{,}179 unique videos, with query types \texttt{v}, \texttt{t}, and \texttt{vt}. For the lightweight subtitle-window baselines, the validation queries are loaded from the official TVR validation metadata, while subtitle windows are built from the official preprocessed subtitle JSONL. Following the repo design, this part of the evaluation is clip-conditioned: the correct \texttt{vid\_name} is known, and the task is to recover the correct temporal window inside that clip.

For the multimodal runs, we also evaluate on the same validation split with the official \xml corpus-level protocol. These runs use the vendored \texttt{third\_party/TVRetrieval} implementation together with the official precomputed query, subtitle, and video features. As a result, the subtitle-window baselines and the \xml family share the same underlying validation queries, but they should be interpreted as two related evaluation regimes: clip-conditioned subtitle retrieval versus corpus-level moment retrieval.

\subsection{Baselines}
We report baselines in the order they build up the system, but the main report emphasis is on the multimodal \xml stage rather than the transcript-only front-end. The first group therefore serves mainly to choose a cheap retrieval module: \bm over subtitle windows, dense subtitle retrieval with MiniLM-L6-v2 sentence embeddings, and sparse+dense fusion. For fusion, we ran the repo's validation sweep over weighted fusion, reciprocal-rank fusion, and \bm-top-$n$ dense reranking, and we use the strongest validation setting, namely weighted fusion with $\alpha = 0.25$. We also keep the oracle-window baseline as an upper bound on what clip-conditioned subtitle windows alone could achieve.

The second group contains the multimodal models that matter most for the systems question in this report. We use the official \xml \texttt{video\_sub + resnet\_i3d} validation result as the full-reference baseline. Our main system then inserts a dense retrieval stage before \xml and only forwards a shortlist of retrieved videos to the expensive multimodal localizer. We evaluate shortlist depths of top-1, top-5, and top-10, with top-5 treated as the primary operating point. We also include the shared-compaction side experiment as a systems ablation because it reduces temporal context inside \xml directly, but it is not our main method. The repo also contains a visual reranking path, but we do not include it in the final validation analysis here because this checkout does not contain the validation-time \clip frame-feature cache required to reproduce that branch locally.

\subsection{Implementation Details}
Subtitle windows are built once offline with a 12\,s window size and a 6\,s stride. During retrieval, the code lazily constructs a per-video \bm index or dense embedding matrix the first time a clip is visited, then reuses that cached state for later queries. Dense query embeddings are encoded in batches on GPU and amortized across the validation sweep. For the fusion runs, the repo measures \bm scoring time, dense scoring time, and fusion time separately, which lets us attribute the cost of each stage instead of only reporting one end-to-end number.

The dense retrieval model is \texttt{all-MiniLM-L6-v2}. The multimodal runs use the upstream \xml implementation together with the official TVR feature release. Validation was executed on a Tesla T4 machine with Python 3.10 and PyTorch 2.5.1. For \xml, we keep the official inference path but run with \texttt{--no\_core\_driver} so that the large HDF5 feature files stay on disk rather than being copied into RAM. This is an infrastructure choice rather than a modeling change, but it matters for making full validation feasible on the available host.

\subsection{Metrics}
For the clip-conditioned subtitle-window baselines, we report mean IoU together with recall at top-1 and top-3 under IoU thresholds 0.3 and 0.5. In the main text we emphasize the IoU 0.5 results because they are the stricter operating point, but the generated summaries retain both thresholds. Latency for these baselines is reported as average online time per query over the full validation sweep with warm caches. The reported numbers include query-time encoding and scoring, but exclude offline subtitle parsing, window construction, and document embedding generation.

For the \xml runs, we use the official TVR metrics: VR for video retrieval, VCMR for corpus-level moment retrieval, and SVMR for within-video localization. Since the official scripts report total validation wall time, we convert runtime to milliseconds per query by dividing the recorded sweep time by the 10{,}895 validation queries. Speedup is always measured relative to the full \xml baseline under the same validation protocol. This makes the main systems question precise: how much corpus-level retrieval and localization quality can be retained while reducing the amount of expensive multimodal inference.
